\uuid{0CXv}
\exo7id{7029}
\titre{exo7 7029}
\auteur{megy}
\organisation{exo7}
\datecreate{2016-10-28}
\isIndication{false}
\isCorrection{false}
\chapitre{Déterminant, système linéaire}
\sousChapitre{Système linéaire, rang}
\cpge{mpsi-mp2i.calcul-algebrique-trigonometrie.b, pcsi.calcul-algebrique-trigonometrie.b, ptsi.calcul-algebrique-trigonometrie.b, mpsi-mp2i.calcul-matriciel-systemes.c, pcsi.calcul-matriciel-systemes-lineaires.c, ptsi.calcul-matriciel-systemes-lineaires.c, bcpst1.systemes-lineaires.a, bcpst1.systemes-lineaires.c}
\module{Algèbre}
\niveau{L2}
\difficulte{}

\contenu{
\question{
Dans une station de métro, les usagers ont à leur disposition un tapis roulant de 300m de long.
Un piéton marchant à vitesse constante fait l'aller-retour sur le tapis roulant. À l'aller, il met 1 minute et
30 secondes. Au retour, à contre-sens, il met 4 minutes et 30 secondes.
Déterminer la vitesse du piéton et celle du tapis roulant.
}
}
